\section{The API}
\label{sec:api}

This is the contract both tiers build against. FastAPI generates the same information as an OpenAPI
schema, served at \code{/docs}, \code{/redoc} and \code{/openapi.json}; the report's API section is
written from that schema so the two cannot differ.

\subsection{Conventions}
\label{sec:apiconventions}

\begin{itemize}
  \item \textbf{Base and format.} Every path starts with \code{/api}. Bodies are JSON in UTF-8,
    except the CSV export (E10) and the progress stream (E18). Names are camelCase, as in
    Assignment 1's contract and \code{predict.py}.
  \item \textbf{Resources.} Paths name things (addresses, transactions, methods, watchlist entries,
    lookups), and the HTTP method says what happens: GET reads and changes nothing; POST computes
    from a body or starts a job; PUT creates or replaces the named entry and can be repeated safely;
    DELETE removes it. Assignment 1's \code{POST /api/lookup} keeps its name because it was promised.
  \item \textbf{Times and amounts.} Times are ISO 8601 in UTC, such as \code{2026-10-20T12:00:00Z}.
    Days are \code{YYYY-MM-DD} in UTC. Amounts in responses are whole satoshis (\code{netSats});
    feature values are in the units people read (BTC, count, hours, ratio). The what-if request
    takes \code{amountBtc}, with up to eight decimal places.
  \item \textbf{Lists.} \code{page} starts at 1; \code{pageSize} defaults to 50 and stops at 200.
    A list answers with its rows, \code{page}, \code{pageSize} and \code{total}. No rows is a 200
    with an empty list, never an error.
  \item \textbf{Errors.} One body for every failure (Section~\ref{sec:errorbody}).
  \item \textbf{Headers.} Every response carries \code{X-Request-ID}. Snapshot data carries an
    \code{ETag} and \code{Cache-Control: no-cache}, so a repeat request costs a 304. Responses over
    1 kB are compressed with gzip.
  \item \textbf{The deciding method.} E3, E5, E6, E7 and E10 accept \code{method=} (S-01). Without
    it they use \code{explorerMethod}, the baseline.
\end{itemize}

\subsection{Every endpoint}
\label{sec:endpoints}

{\small
\begin{longtable}{@{}L{0.7cm}L{1.3cm}L{6.6cm}L{4.0cm}L{1.7cm}@{}}
\toprule
\textbf{No.} & \textbf{Method} & \textbf{Path} & \textbf{Returns} & \textbf{Tier} \\
\midrule
\endhead
E1 & GET & \code{/api/health} & Ready or starting; what is loaded & \must \\
E2 & GET & \code{/api/config} & Snapshot, limits, thresholds, feature names, capabilities & \must \\
E3 & GET & \code{/api/addresses} & The addresses that can be opened, with counts & \must{} (A1) \\
E4 & POST & \code{/api/lookup} & Whether a typed address is supported & \must{} (A1) \\
E5 & GET & \code{/api/addresses/{address}/timeline} & Per-day counts, totals, range, history status & \must{} (A1) \\
E6 & GET & \code{/api/addresses/{address}/transactions} & A filtered, sorted page of transaction rows & \must{} (A1, extended) \\
E7 & GET & \code{/api/transactions/{txid}?address=} & One transaction: facts, features, reason, every model's verdict & \must{} (A1) \\
E8 & POST & \code{/api/addresses/{address}/what-if} & A hypothetical transaction's verdict from every model & \must \\
E9 & POST & \code{/api/addresses/{address}/what-if/curve} & Scores over a range of amounts & \must \\
E10 & GET & \code{/api/addresses/{address}/export} & The scored transactions as a CSV or JSON file & \must \\
E11 & GET & \code{/api/methods} & The nine methods, their thresholds and test results & \must \\
E12 & POST & \code{/api/explanations} & The Assistant's guarded explanation, or the template & \should{} S-05 \\
E13 & GET & \code{/api/watchlist} & The watched addresses with labels and notes & \should{} S-04 \\
E14 & PUT & \code{/api/watchlist/{address}} & Adds or replaces an entry & \should{} S-04 \\
E15 & DELETE & \code{/api/watchlist/{address}} & Removes an entry & \should{} S-04 \\
E16 & POST & \code{/api/live/lookups} & Starts a live lookup; answers 202 with its identifier & \should{} S-09 \\
E17 & GET & \code{/api/live/lookups/{id}} & The lookup's status and progress & \should{} S-09 \\
E18 & GET & \code{/api/live/lookups/{id}/events} & The same as a stream of events & \should{} S-09 \\
E19 & DELETE & \code{/api/live/lookups/{id}} & Cancels the lookup & \should{} S-09 \\
E20 & DELETE & \code{/api/addresses/{address}} & Forgets a fetched address & \should{} S-09 \\
E21 & GET & \code{/api/addresses/{address}/surprisal} & TxGPT's surprisal per field & \could{} C-01 \\
E22 & GET & \code{/api/addresses/{address}/profile} & The address's activity profile & \could{} C-02 \\
\bottomrule
\end{longtable}
}

The Must tier uses GET and POST across eleven endpoints, which meets the brief's ``at least two
HTTP methods''. S-04 adds PUT and DELETE on a resource where they mean what REST says they mean.

\subsection{The error body and the status codes}
\label{sec:errorbody}

\begin{lstlisting}[style=json]
{
  "status": "invalid",
  "message": "The amount must be greater than zero.",
  "hint": "Enter an amount in BTC, such as 0.05.",
  "fields": [ { "name": "amountBtc", "message": "Must be greater than 0." } ],
  "requestId": "0f3c9a1e7b2d"
}
\end{lstlisting}

\code{status} is for the code, \code{message} and \code{hint} are plain words that are safe to
show, \code{fields} names the form fields at fault, and \code{requestId} matches the server's log
line. The first three are Assignment 1's bodies unchanged; \code{fields} and \code{requestId} are
additions. Some statuses add their own keys, as Assignment 1 planned: \code{examples} with
\code{unsupported}, and \code{transaction}, \code{history} and \code{firstScoredDate} with
\code{insufficient_history}.

{\small
\begin{longtable}{@{}L{0.9cm}L{3.6cm}L{5.2cm}L{4.2cm}@{}}
\toprule
\textbf{Code} & \textbf{\code{status}} & \textbf{When} & \textbf{Raised by} \\
\midrule
\endhead
200 & -- & Success with a body & \\
201 & -- & A watchlist entry was created & E14 \\
202 & -- & A live lookup was accepted and is running & E16 \\
204 & -- & Removed; no body & E15, E19, E20 \\
304 & -- & The cached copy is still right & snapshot GETs \\
400 & \code{invalid} & A malformed address, date, transaction ID, number, choice or body; an unknown field; \code{from} after \code{to} & Pydantic validation, the address checker, \code{ValueError} \\
403 & \code{disabled} & Live lookup is switched off on this server & E16 to E20 \\
404 & \code{unsupported} & A valid address that is not in the snapshot or the live cache; carries three examples & \code{Unsupported} \\
404 & \code{not_found} & An unknown transaction for that address, watchlist entry, lookup or path & \code{Unsupported}, the repositories \\
409 & \code{conflict} & The watchlist is full; a lookup for that address is already running; a snapshot address cannot be forgotten & E14, E16, E20 \\
413 & \code{too_large} & A body over the size limit & middleware \\
415 & \code{invalid} & A body that is not JSON & FastAPI \\
422 & \code{insufficient_history} & Too little earlier history to score: the address, the transaction, or the what-if's time & \code{InsufficientHistory} \\
429 & \code{rate_limited} & Too many live lookups, or the Assistant is busy; carries \code{Retry-After} & E12, E16 \\
500 & \code{server_error} & Anything unexpected. A fixed message and the request identifier; the trace goes to the log only & the catch-all handler \\
502 & \code{upstream_error} & Esplora answered with an error & E16 to E18 \\
503 & \code{unavailable} & The models are still loading, or a file is missing & E1; \code{FileNotFoundError} \\
504 & \code{upstream_timeout} & Esplora did not answer in time & E16 to E18 \\
\bottomrule
\end{longtable}
}

\begin{caution}
FastAPI's own validation error is a 422 with a \code{detail} list. The handler for
\code{RequestValidationError} is replaced so that it answers 400 with the body above (D-A3-5). A
test posts a wrong type to every POST and PUT endpoint and checks for 400 and \code{fields}.
\end{caution}

\subsection{Endpoints in detail}

Shapes marked A1 are those of \fp{frontend/src/api/types.ts} in the prototype. Full example bodies
are in Appendix~\ref{app:json}.

\subsubsection*{E1 \quad GET /api/health}
Answers 200 with \code{status: "ok"}, the application version, \code{snapshotDate},
\code{addressCount}, \code{methodCount}, \code{assistant} (\code{ready}, \code{not_installed},
\code{loading} or \code{off}), \code{liveLookup} (true or false) and \code{uptimeSeconds}. While the
models load it answers 503 with \code{status: "unavailable"}. Never cached. The front end polls it
for the service-status indicator, and the README uses it as the first check.

\subsubsection*{E2 \quad GET /api/config}
Everything the front end would otherwise hard-code: \code{snapshot} (date, source, link, address
count, method), \code{limitationLine}, \code{features} (identifier, name, description and unit of F1
to F4), \code{history} (the minimum of 10 transactions over 14 days), \code{thresholds} (the flag
line 5.0 and the contributing line 2.0), \code{defaultMethod}, the \code{web} block of
Section~\ref{sec:config}, \code{whatIf.defaultTime} (noon UTC on the snapshot date) and
\code{capabilities} (\code{assistant}, \code{liveLookup}).

\subsubsection*{E3 \quad GET /api/addresses}
A1's \code{AddressesResponse}: \code{snapshot} and \code{addresses}. Each address has
\code{address}, \code{firstSeen}, \code{lastSeen}, \code{txCount}, \code{scoredCount},
\code{flaggedCount} and \code{history} (\code{sufficient} or \code{insufficient}). Added:
\code{kind} (the address type, such as P2PKH), \code{source} (\code{snapshot} or \code{live}) and,
with S-04, \code{watch} (label and note, or null). Main addresses come first, busiest first; the
short-history ones last, as in the terminal Explorer.

\subsubsection*{E4 \quad POST /api/lookup}
Body: \code{address}. 200 with \code{status: "supported"}, the normalised \code{address} and its
\code{summary}. 400 \code{invalid} with the checker's specific message, such as ``This address fails
its checksum, so it probably contains a typo.'' 404 \code{unsupported} with the coverage sentence
and three \code{examples}; with S-09 it adds \code{canLookupLive}.

\subsubsection*{E5 \quad GET /api/addresses/\{address\}/timeline}
A1's \code{TimelineResponse}: \code{address}, \code{snapshot}, \code{range} (first and last
transaction time), \code{days} (only days with transactions, oldest first, each with
\code{txCount} and \code{flaggedCount}), \code{totals} (\code{txCount}, \code{scoredCount},
\code{flaggedCount}, \code{flaggedDays}) and \code{history} (\code{status}, \code{txCount},
\code{spanDays}, \code{required}, \code{firstScoredDate}). A short-history address answers 200 with
\code{history.status: "insufficient"} and \code{scoredCount: 0}, so the screen can list its
transactions unscored. Errors: 400, 404 \code{unsupported}.

\subsubsection*{E6 \quad GET /api/addresses/\{address\}/transactions}
\begin{center}
\small
\begin{tabularx}{\textwidth}{@{}L{2.4cm}L{4.3cm}Y@{}}
\toprule
\textbf{Query} & \textbf{Values} & \textbf{Meaning} \\
\midrule
\code{date} & \code{YYYY-MM-DD} & One UTC day (A1's parameter) \\
\code{from}, \code{to} & \code{YYYY-MM-DD} & An inclusive range of days; not with \code{date} \\
\code{status} & \code{flagged}, \code{not-flagged}, \code{not-scored}; may repeat & Keep these statuses \\
\code{direction} & \code{in}, \code{out}, \code{none}; may repeat & Keep these directions \\
\code{sort} & \code{time}, \code{amount}, \code{score}; a leading minus reverses & Default \code{time} \\
\code{page}, \code{pageSize} & whole numbers & Default 1 and 50; at most 200 a page \\
\code{method} & a method identifier & Whose flags and scores to show (S-01) \\
\bottomrule
\end{tabularx}
\end{center}
Answers \code{address}, \code{date} (or null), \code{transactions}, \code{page}, \code{pageSize}
and \code{total}. A row is A1's \code{TransactionRow} (\code{txid}, \code{time}, \code{direction},
\code{netSats}, \code{status}) plus \code{score} (the deciding method's, or null),
\code{mainReason} (the template's strongest reason, such as ``Longer gap than usual''),
\code{mainZ} and \code{flaggedBy} (how many methods flag it). A day or a filter that matches
nothing is a 200 with an empty list. Errors: 400 for a malformed or contradictory query, 404
\code{unsupported}.

The scatter chart asks for up to 200 rows a page and joins the pages; the largest snapshot address
needs four requests.

\subsubsection*{E7 \quad GET /api/transactions/\{txid\}?address=}
A1's \code{TransactionDetail}: the facts (\code{txid}, \code{address}, \code{time},
\code{blockHeight}, \code{direction}, \code{receivedSats}, \code{sentSats}, \code{netSats},
\code{sourceUrl}), \code{scored}, \code{flagged}, \code{score}, \code{method}, \code{features}
(for each of F1 to F4: \code{value}, \code{usual}, \code{z} and \code{role}), \code{comparison}
(what each feature was compared with: how many earlier transactions, from when to when) and
\code{limitation}. Added, all already returned by \code{predict.py}: \code{explanation} (the
template text), \code{scores} and \code{flags} for every method, and \code{source}. Errors: 400
(a transaction ID that is not 64 hexadecimal characters, or no \code{address}); 404
\code{unsupported} (the address); 404 \code{not_found} (the transaction is not that address's);
422 \code{insufficient_history} with the transaction's facts when it could not be scored.

\subsubsection*{E8 \quad POST /api/addresses/\{address\}/what-if}
\label{sec:e8}
\begin{center}
\small
\begin{tabularx}{\textwidth}{@{}L{2.2cm}L{3.1cm}Y@{}}
\toprule
\textbf{Field} & \textbf{Type} & \textbf{Rule} \\
\midrule
\code{amountBtc} & number or decimal string & Greater than 0; at least 0.00000001; at most 21,000,000; at most 8 decimal places \\
\code{direction} & \code{"in"} or \code{"out"} & Required \\
\code{time} & ISO 8601 with a zone & Not before the address's first transaction (else 422); not more than \code{maxGapYears} after its last (else 400) \\
\bottomrule
\end{tabularx}
\end{center}
200: a transaction detail as E7 with \code{txid: "what-if"}, \code{sourceUrl: null}, the template
\code{explanation}, \code{scores} and \code{flags} from the eight methods that can score new input,
\code{notAvailable: ["txgpt_method"]} and a \code{note} that the transaction is hypothetical.
Nothing is stored. Errors: 400 \code{invalid} with \code{fields}; 404 \code{unsupported}; 422
\code{insufficient_history} when fewer than 10 transactions or 14 days of history precede the
chosen time, with the counts and the first date that can be scored.

\subsubsection*{E9 \quad POST /api/addresses/\{address\}/what-if/curve}
Body: \code{direction} and \code{time} as E8, then either \code{amountsBtc} (a list of up to 61
amounts) or \code{minBtc}, \code{maxBtc} and \code{points}. With neither, the server spaces 41
amounts evenly on a log scale from one hundredth of the address's usual amount in that direction to
a thousand times it. 200: \code{amountsBtc}, \code{usualBtc}, \code{decidingMethod}, and for each
method its \code{threshold}, \code{scores} and \code{flags}, one per amount. One feature build
scores every amount (Section~\ref{sec:whatif}). Errors as E8.

\subsubsection*{E10 \quad GET /api/addresses/\{address\}/export}
Query: \code{format} (\code{csv}, the default, or \code{json}) and E6's filters. Answers a file
with \code{Content-Disposition: attachment}, named from the first eight characters of the address
and the snapshot date. CSV columns: address, transaction ID, time (UTC), direction, received, sent
and net amount in BTC, scored, flagged, deciding method, score, each feature's value, usual value
and z, the main reason, the methods that flag it, the source link, the snapshot date and, on
flagged rows, the limitation line. A cell that starts with \code{=}, \code{+}, \code{-} or
\code{@} is written with a leading apostrophe, so a spreadsheet does not run it as a formula.

\subsubsection*{E11 \quad GET /api/methods}
\code{defaultMethod} and \code{methods}. Each has \code{id}, \code{name}, \code{family} (statistics,
outlier detection, clustering, regression, combination, language model), a one-sentence
\code{description}, its \code{threshold}, \code{scoresNewInput} and \code{evaluation}: mean recall
on planted transactions with its 95\% interval, recall per pattern (spike, burst, dormant) and the
flag rate on real transactions. The evaluation numbers are read from copies of Assignment 2's
\fp{eval_significance.csv} and \fp{eval_selection.csv}; none is typed.

\subsubsection*{E12 \quad POST /api/explanations \quad (S-05)}
Body: \code{address} and either \code{txid} or a \code{whatIf} object as E8. 200: \code{txid},
\code{flagged}, \code{explanation}, \code{source} (\code{text150m} or \code{template}),
\code{cached}, \code{attempts} (each writer tried and the guardrail checks it failed),
\code{skippedWriters} (a writer this machine cannot run, and why), \code{limitation} and
\code{seconds}. A rejected answer is never returned. Errors: 400, 404, 422; 429 while another
answer is being written.

\subsubsection*{E13 to E15 \quad the watchlist \quad (S-04)}
\code{GET /api/watchlist} lists entries (\code{address}, \code{label}, \code{note}, \code{addedAt},
\code{updatedAt} and the address's summary). \code{PUT /api/watchlist/{address}} with
\code{label} (40 characters at most) and \code{note} (500 at most) answers 201 for a new entry and
200 for a replaced one; 404 \code{unsupported} for an address that cannot be opened; 409 when the
list holds 200 entries. \code{DELETE} answers 204, or 404 \code{not_found}.

\subsubsection*{E16 to E20 \quad live lookup \quad (S-09, CR-08)}
\code{POST /api/live/lookups} with \code{address} and optional \code{maxTx} (50 to 1,000; default
500) answers 202 with the lookup's \code{id}, its \code{status} and links to E17 and E18, and a
\code{Location} header. A fresh cached result answers 200 at once. \code{GET} E17 returns
\code{status} (\code{queued}, \code{fetching}, \code{scoring}, \code{done}, \code{failed},
\code{cancelled}), \code{progress} (pages and transactions fetched, the cap) and, on failure, an
error body. E18 sends the same as server-sent events named \code{progress}, \code{status},
\code{done} and \code{failed}, with a comment line every 15 seconds to keep the connection open.
E19 cancels. Once done, the address answers E5 to E10 like a snapshot address, with
\code{source: "live"} and the time it was fetched; E20 forgets it.

\subsubsection*{E21, E22 \quad (Could)}
E21 returns, for a snapshot address, each transaction's surprisal in total and for the six token
fields (gap, direction, amount, inputs, outputs, fee). E22 returns the address's profile from
Assignment 2's profile model: its label, its eight address-level features and the profile's medians.

\subsection{What the report documents}

The brief asks for at least four endpoints. The report gives the table of
Section~\ref{sec:endpoints}, the error body, and four endpoints in full with a request and a real
response each: E4 (a POST with all three outcomes), E5, E7 and E8. The complete reference, exported
from \code{/openapi.json}, goes in the appendix, which does not count towards the eight pages.

\begin{done}
Every endpoint in a finished tier has a response model, appears in \code{/docs} with a summary and
an example, and has one test for its success and one for each error in its row. \code{GET} E3 to E7
return exactly what \code{Predictor} returns for all 28 addresses (the parity test of
Section~\ref{sec:tests}).
\end{done}
